\chapter{Typy i przestrzenie typów Stone'a}
\label{ch:mod-types}
\section{Od pojedynczej formuły do opisu elementu}
Formuła $x>0$ wyróżnia pewną część modelu, lecz nie opisuje
wszystkich własności żadnego konkretnego elementu. Możemy jednocześnie
żądać spełniania wielu formuł, nawet nieskończenie wielu. Typ
jest takim zbiorem wymagań, zgodnym z ustaloną teorią.
Typ zupełny rozstrzyga każdą formułę dotyczącą rozważanej krotki.
Nie należy go utożsamiać z elementem: różne elementy mogą
mieć ten sam typ, a typ może nie być realizowany w danym modelu.

Pracujemy w ZFC. Ustalamy sygnaturę $\mathcal L$ i liczbę
naturalną $n\geq1$. Przez $\bar x=(x_1,\ldots,x_n)$ rozumiemy
listę parami różnych zmiennych, a przez $F_n(\mathcal L)$ zbiór
formuł, których zmienne wolne należą do tej listy.
Wykorzystujemy semantykę z definicji \ref{def:sem-satisfaction}
oraz twierdzenie o zwartości \ref{thm:compactness}.

\section{Typy nad teorią}
\begin{definicja}[Typ częściowy, realizacja i pomijanie]
\label{def:mod-partial-type}
Niech $T$ będzie spełnialną teorią języka $\mathcal L$.
Zbiór $p\subseteq F_n(\mathcal L)$ nazywamy typem częściowym
nad $T$, jeżeli istnieją model $\mathfrak M\models T$
i krotka $\bar a\in M^n$ takie, że
$\mathfrak M\models\varphi(\bar a)$ dla każdej $\varphi\in p$.
Mówimy wtedy, że $\bar a$ realizuje $p$ w $\mathfrak M$.
Model $\mathfrak M$ pomija $p$, gdy żadna krotka z $M^n$
nie realizuje $p$.
\end{definicja}

\begin{lemat}[Kryterium skończonej spełnialności typu]
\label{lem:mod-type-compactness}
Zbiór $p\subseteq F_n(\mathcal L)$ jest typem częściowym nad $T$
wtedy i tylko wtedy, gdy każdy skończony podzbiór $p_0\subseteq p$
jest realizowany w pewnym modelu $T$.
\end{lemat}
\begin{proof}
Konieczność wynika z definicji. Dla wystarczalności dodajemy do
$\mathcal L$ nowe stałe $\bar d=(d_1,\ldots,d_n)$.
Teoria $T\cup\{\varphi(\bar d):\varphi\in p\}$ w rozszerzonym
języku jest skończenie spełnialna: dla danego skończonego
fragmentu wybieramy model $T$ i krotkę realizującą odpowiadający
skończony podzbiór $p$, po czym interpretujemy nowe stałe
przez współrzędne tej krotki. Podstawienie stałych za zmienne
wolne jest dopuszczalne i zachowuje spełnianie zgodnie
z lematem \ref{thm:sem-substitution}. Ze zwartości istnieje
model całej teorii rozszerzonej. Jego redukt i wartości stałych
realizują $p$.
\end{proof}

\begin{definicja}[Typ zupełny]
\label{def:mod-complete-type}
Typ częściowy $p$ nad $T$ nazywamy zupełnym, jeżeli dla każdej
$\varphi\in F_n(\mathcal L)$ dokładnie jedna z formuł
$\varphi,\neg\varphi$ należy do $p$.
Zbiór wszystkich takich typów oznaczamy $S_n(T)$.
Dla $\mathfrak M\models T$ i $\bar a\in M^n$ kładziemy
\[
\operatorname{tp}^{\mathfrak M}(\bar a)
=\{\varphi\in F_n(\mathcal L):\mathfrak M\models\varphi(\bar a)\}.
\]
\end{definicja}

\begin{twierdzenie}[Maksymalność i uzupełnianie typów]
\label{thm:mod-type-completion}
Typ częściowy jest zupełny wtedy i tylko wtedy, gdy jest
maksymalny względem zawierania wśród typów częściowych nad $T$.
Każdy typ częściowy zawiera się w typie zupełnym, a każdy
typ zupełny jest typem pewnej krotki w pewnym modelu $T$.
\end{twierdzenie}
\begin{proof}
Jeżeli $p$ jest typem częściowym, wybieramy jego realizację
$\bar a$ w modelu $\mathfrak M\models T$. Typ
$\operatorname{tp}^{\mathfrak M}(\bar a)$ zawiera $p$
i jest zupełny, ponieważ semantyka klasyczna rozstrzyga
prawdziwość każdej formuły przy ustalonej krotce.
Jeżeli $p$ jest maksymalny, musi być równy temu typowi.
Jeżeli natomiast $p$ jest zupełny i $\varphi\notin p$,
to $\neg\varphi\in p$, a zatem dodanie $\varphi$
uniemożliwia realizację. Dowodzi to maksymalności, a także
równości $p=\operatorname{tp}^{\mathfrak M}(\bar a)$ dla każdej
realizacji $p$.
\end{proof}

\section{Parametry i rozszerzenia elementarne}
\begin{definicja}[Typ nad zbiorem parametrów]
\label{def:mod-type-over-parameters}
Niech $\mathfrak M$ będzie strukturą języka $\mathcal L$
i $A\subseteq M$. Przez $\mathcal L(A)$ oznaczamy język
uzyskany przez dodanie nowej stałej $c_a$ dla każdego $a\in A$.
Niech $\mathfrak M_A$ będzie rozwinięciem $\mathfrak M$
interpretującym $c_a$ przez $a$, a
$T_A=\operatorname{Th}(\mathfrak M_A)$ zbiorem wszystkich
zdań prawdziwych w $\mathfrak M_A$.
Typem częściowym (odpowiednio zupełnym) nad $A$ nazywamy
typ częściowy (odpowiednio zupełny) nad $T_A$.
Piszemy $S_n(A)=S_n(T_A)$, przy ustalonym modelu $\mathfrak M$.
Dla $\bar b\in M^n$ piszemy
$\operatorname{tp}^{\mathfrak M}(\bar b/A)
=\operatorname{tp}^{\mathfrak M_A}(\bar b)$.
\end{definicja}

Zapis $\varphi(\bar x,\bar a)$ oznacza formułę języka $\mathcal L(A)$
z nazwami skończonej krotki parametrów $\bar a$.
Zbiór parametrów nie musi być uniwersum podstruktury.

\begin{twierdzenie}[Realizacja w rozszerzeniu elementarnym]
\label{thm:mod-type-elementary-extension}
Dla $p\subseteq F_n(\mathcal L(A))$ następujące warunki
są równoważne:
\begin{enumerate}[label=(\roman*)]
\item $p$ jest typem częściowym nad $A$;
\item każdy skończony podzbiór $p$ jest realizowany w $\mathfrak M_A$;
\item istnieje $\mathfrak N\succeq\mathfrak M$ oraz
$\bar b\in N^n$ realizująca $p$ przy nazwach parametrów z $A$.
\end{enumerate}
\end{twierdzenie}
\begin{proof}
Z (i) wynika (ii): dla skończonego $p_0\subseteq p$ zdanie
$\exists x_1\cdots\exists x_n\,\bigwedge_{\varphi\in p_0}\varphi$
jest prawdziwe w pewnym modelu $T_A$. Ponieważ $T_A$ jest
pełną teorią struktury $\mathfrak M_A$, zdanie to jest prawdziwe
również w $\mathfrak M_A$ (w przeciwnym razie jego negacja
należałaby do $T_A$). Dla $p_0=\varnothing$ przyjmujemy
koniunkcję równą $x_1=x_1$.

Aby uzyskać (iii) z (ii), dodajemy nazwy wszystkich elementów
$M$ oraz nowe stałe $\bar d$. Rozważamy teorię
\[
\operatorname{ED}(\mathfrak M)\cup
\{\varphi(\bar d):\varphi\in p\},
\]
gdzie diagram elementarny pochodzi z definicji
\ref{def:ls-elementary-diagram}, a nazwy z $A$ są wspólne
dla obu języków. Każdy skończony fragment ma model:
rozszerzamy $\mathfrak M_M$ o interpretację $\bar d$
przez realizację odpowiedniego skończonego podzbioru $p$.
Z twierdzenia o zwartości istnieje model całej teorii.
Lemat \ref{lem:ls-diagram-embedding} daje zanurzenie
elementarne $\mathfrak M$ w jego redukt do $\mathcal L$.
Przenosimy strukturę przez bijekcję, która utożsamia obraz
każdego $a\in M$ z $a$, a pozostałe elementy z nowymi elementami
spoza $M$. Otrzymujemy dosłowne rozszerzenie elementarne
i krotkę realizującą $p$.

Wreszcie (iii) implikuje (i), ponieważ rozszerzenie elementarne
po nazwaniu elementów $A$ spełnia wszystkie zdania $T_A$.
\end{proof}

\begin{przyklad}[Typ niestandardowego elementu]
\label{ex:mod-type-nonstandard}
W standardowej arytmetyce $\mathfrak N$, w języku zawierającym
porządek, rozważmy nad $\N$ zbiór formuł
$p(x)=\{c_k<x:k\in\N\}$.
Każdy skończony podzbiór jest realizowany przez liczbę większą
od wszystkich występujących w nim parametrów. Cały typ nie
jest realizowany w $\mathfrak N$, ponieważ dla dowolnego $m$
formuła $c_m<x$ nie jest spełniona przy $x=m$.
Z twierdzenia \ref{thm:mod-type-elementary-extension}
typ ma realizację w pewnym rozszerzeniu elementarnym.
Konkretną realizację daje ultrapotęga z przykładu
\ref{ex:mod-ultrapower-naturals}.
\end{przyklad}

\section{Topologia przestrzeni typów}
\begin{definicja}[Zbiory bazowe i topologia Stone'a]
\label{def:mod-stone-topology}
Dla $\varphi\in F_n(\mathcal L)$ definiujemy
\[
[\varphi]=\{p\in S_n(T):\varphi\in p\}.
\]
Topologią Stone'a na $S_n(T)$ nazywamy topologię
o bazie wszystkich zbiorów $[\varphi]$.
\end{definicja}

\begin{lemat}[Rachunek zbiorów bazowych]
\label{lem:mod-stone-basic}
Dla dowolnych $\varphi,\psi\in F_n(\mathcal L)$ zachodzą
\[
\begin{aligned}
[\varphi\land\psi]&=[\varphi]\cap[\psi],\\
[\varphi\lor\psi]&=[\varphi]\cup[\psi],\\
[\neg\varphi]&=S_n(T)\setminus[\varphi],\\
[x_1=x_1]&=S_n(T),\qquad [\neg(x_1=x_1)]=\varnothing.
\end{aligned}
\]
W szczególności wskazana rodzina rzeczywiście jest bazą
topologii, a każdy zbiór bazowy jest otwarty i domknięty.
\end{lemat}
\begin{proof}
Każdy typ zupełny jest typem krotki w modelu $T$
(twierdzenie \ref{thm:mod-type-completion}). Równości
wynikają więc bezpośrednio z semantyki spójników i równości.
Rodzina pokrywa całą przestrzeń i jest zamknięta na skończone
przecięcia, co zapewnia warunki bazy.
\end{proof}

\begin{twierdzenie}[Zwartość przestrzeni typów]
\label{thm:mod-stone-compact}
Przestrzeń $S_n(T)$ jest niepusta, zwarta, Hausdorffa
i zerowymiarowa, tj. ma bazę złożoną ze zbiorów otwartych
i domkniętych.
\end{twierdzenie}
\begin{proof}
Niepustość wynika z istnienia modelu $T$ i jego niepustego
uniwersum. Jeśli $p\ne q$, pewna formuła należy do jednego
typu, a nie należy do drugiego; po ewentualnej zamianie
typów mamy $\varphi\in p$, $\neg\varphi\in q$.
Rozłączne zbiory $[\varphi]$ i $[\neg\varphi]$ rozdzielają
$p$ i $q$. Zerowymiarowość wynika z lematu \ref{lem:mod-stone-basic}.

Dla zwartości wystarczy sprawdzić pokrycia zbiorami bazowymi:
każde pokrycie otwarte można uszczegółowić do pokrycia bazowego,
a skończone podpokrycie uszczegółowienia daje skończone
podpokrycie pierwotnego. Załóżmy, że rodzina
$([\varphi_i])_{i\in I}$ pokrywa $S_n(T)$, lecz żadna skończona
podrodzina nie pokrywa przestrzeni. Dla każdego skończonego
$J\subseteq I$ istnieje typ zupełny zawierający wszystkie
$\neg\varphi_j$, $j\in J$. Zatem
$\{\neg\varphi_i:i\in I\}$ jest skończenie realizowalny
w modelach $T$. Lemat \ref{lem:mod-type-compactness} i twierdzenie
\ref{thm:mod-type-completion} dają typ zupełny zawierający
wszystkie te negacje. Nie należy on do żadnego zbioru
pokrycia, co jest sprzecznością.
\end{proof}

\begin{wniosek}[Domknięte zbiory i typy częściowe]
\label{cor:mod-stone-closed-types}
Dla $p\subseteq F_n(\mathcal L)$ zbiór
$X_p=\{q\in S_n(T):p\subseteq q\}$ jest domknięty.
Każdy zbiór domknięty w $S_n(T)$ ma tę postać.
Ponadto $X_p\ne\varnothing$ dokładnie wtedy, gdy $p$
jest typem częściowym.
\end{wniosek}
\begin{proof}
Mamy $X_p=\bigcap_{\varphi\in p}[\varphi]$. Jeżeli $C$
jest domknięty, jego dopełnienie jest sumą zbiorów bazowych,
powiedzmy $\bigcup_{i\in I}[\varphi_i]$, więc
$C=X_{\{\neg\varphi_i:i\in I\}}$ na mocy
lematu \ref{lem:mod-stone-basic}. Ostatnia równoważność
wynika z twierdzenia \ref{thm:mod-type-completion}.
\end{proof}

\begin{wniosek}[Zbiory otwarto-domknięte]
\label{cor:mod-stone-clopens}
Każdy zbiór otwarty i domknięty w $S_n(T)$ ma postać
$[\varphi]$ dla pewnej formuły $\varphi$.
\end{wniosek}
\begin{proof}
Taki zbiór jest domkniętym podzbiorem przestrzeni zwartej,
a więc jest zwarty: do dowolnego jego pokrycia otwartego
w przestrzeni dodajemy jego otwarte dopełnienie i stosujemy
zwartość całej przestrzeni. Ponieważ zbiór jest otwarty,
pokrywają go zbiory bazowe zawarte w nim. Skończone
podpokrycie daje formułę będącą alternatywą odpowiednich
formuł. Dla zbioru pustego wybieramy $\neg(x_1=x_1)$.
\end{proof}

\section{Algebra Boole'a formuł i jej ultrafiltry}
\begin{definicja}[Algebra formuł modulo teorię]
\label{def:mod-lindenbaum-algebra}
Na $F_n(\mathcal L)$ definiujemy
$\varphi\equiv_T\psi$ wtedy i tylko wtedy, gdy
\[
T\models\forall x_1\cdots\forall x_n\,(\varphi\leftrightarrow\psi).
\]
Przez $B_n(T)$ oznaczamy zbiór klas równoważności, a klasę
formuły $\varphi$ przez $\langle\varphi\rangle_T$.
Działania koniunkcji, alternatywy i negacji na klasach
pochodzą od odpowiednich spójników; elementami $0$ i $1$
są klasy sprzeczności i tautologii.
\end{definicja}

\begin{twierdzenie}[Reprezentacja Stone'a dla algebry formuł]
\label{thm:mod-stone-boolean}
$B_n(T)$ jest algebrą Boole'a, a odwzorowanie
$\langle\varphi\rangle_T\mapsto[\varphi]$ jest izomorfizmem
z algebrą zbiorów otwartych i domkniętych w $S_n(T)$.
Punkty $S_n(T)$ odpowiadają bijektywnie ultrafiltrom
$B_n(T)$ przez przyporządkowanie
\[
p\longmapsto U_p=\{\langle\varphi\rangle_T:\varphi\in p\}.
\]
Bijekcja jest homeomorfizmem, gdy zbiór ultrafiltrów
wyposażymy w topologię o bazie
$\widehat b=\{U:b\in U\}$ dla $b\in B_n(T)$.
\end{twierdzenie}
\begin{proof}
Równoważność semantyczna jest relacją równoważności
zachowywaną przez spójniki. Jeżeli $\varphi\equiv_T\psi$,
to każda krotka w każdym modelu $T$ spełnia je jednocześnie
albo nie spełnia żadnej. Z twierdzenia
\ref{thm:mod-type-completion} wynika $[\varphi]=[\psi]$.
Odwrotnie, jeśli formuły nie są równoważne modulo $T$,
pewna krotka w modelu $T$ odróżnia ich prawdziwość,
a jej typ należy dokładnie do jednego ze zbiorów
$[\varphi]$, $[\psi]$. Otrzymujemy zatem różnowartościowe
odwzorowanie klas. Zachowuje ono działania na mocy lematu
\ref{lem:mod-stone-basic} i jest na na mocy wniosku
\ref{cor:mod-stone-clopens}. To dowodzi również aksjomatów
algebry Boole'a, ponieważ rodzina zbiorów otwartych
i domkniętych jest zamknięta na działania zbiorowe.

Dla typu zupełnego $p$ wybieramy jego realizację.
Prawdziwe przy niej formuły tworzą właściwy filtr klas:
tautologia jest prawdziwa, sprzeczność fałszywa, koniunkcja
prawdziwych formuł prawdziwa, a implikacja semantyczna
zachowuje prawdziwość. Zupełność typu daje wybór dokładnie
jednego z $b,\neg b$, czyli ultrafiltr (kryterium
\ref{lem:bpi-ultrafilter-criterion}).

Dla ultrafiltru $U$ kładziemy
$p_U=\{\varphi:\langle\varphi\rangle_T\in U\}$.
Skończona koniunkcja formuł należących do $p_U$
ma klasę w $U$, zatem różną od zera. Nie jest więc
sprzeczna modulo $T$, czyli pewna krotka w pewnym modelu
$T$ ją spełnia. Lemat \ref{lem:mod-type-compactness}
daje typ częściowy; własność ultrafiltru daje zupełność.
Obie konstrukcje są wzajemnie odwrotne (przynależność
formuły do typu zupełnego zależy tylko od jej klasy).
Zbiory $[\varphi]$ odpowiadają dokładnie zbiorom
$\widehat{\langle\varphi\rangle_T}$, które tworzą bazę
na zbiorze ultrafiltrów, gdyż $\widehat 1$ jest całą
przestrzenią, a $\widehat b\cap\widehat c=\widehat{b\land c}$.
Wynika stąd homeomorfizm.
\end{proof}

\section{Typy izolowane i przykłady}
\begin{definicja}[Typ izolowany]
\label{def:mod-isolated-type}
Typ $p\in S_n(T)$ jest izolowany, jeśli istnieje
$\varphi\in p$ taka, że $[\varphi]=\{p\}$.
\end{definicja}

\begin{twierdzenie}[Izolacja a konsekwencja semantyczna]
\label{thm:mod-isolated-type-characterization}
Dla $p\in S_n(T)$ i $\varphi\in p$ warunek
$[\varphi]=\{p\}$ jest równoważny temu, że
dla każdej $\psi\in p$ zachodzi
\[
T\models\forall\bar x\,(\varphi\to\psi).
\]
Jeżeli ponadto $T$ jest teorią zupełną, każdy typ izolowany
jest realizowany w każdym modelu $T$.
\end{twierdzenie}
\begin{proof}
Jeżeli implikacja nie zachodzi, istnieje krotka w modelu
$T$ spełniająca $\varphi$ i $\neg\psi$. Jej typ jest
innym punktem $[\varphi]$. Odwrotnie, jeśli wszystkie
implikacje zachodzą, każdy typ zawierający $\varphi$
zawiera $p$ (stosujemy jego realizację), a więc jest równy
$p$ przez maksymalność. Dla zupełnego $T$ zdanie
$\exists\bar x\,\varphi$ należy do konsekwencji $T$,
ponieważ zachodzi w pewnym modelu $T$ i wszystkie modele
$T$ są elementarnie równoważne. Każda jego realizacja
ma typ $p$.
\end{proof}

\begin{przyklad}[Nieskończony zbiór w języku równości]
\label{ex:mod-type-pure-equality}
Niech $M$ będzie nieskończonym zbiorem w języku samej
równości, a $A\subseteq M$. Typy $S_1(A)$ to typy
$p_a=\operatorname{tp}^{\mathfrak M}(a/A)$ dla $a\in A$
oraz dokładnie jeden typ $p_\infty$ zawierający
$\{x\ne c_a:a\in A\}$.
Istnienie ostatniego typu wynika z twierdzeń
\ref{thm:mod-type-elementary-extension} i
\ref{thm:mod-type-completion}, gdyż skończenie wiele
nierówności można zrealizować w nieskończonym $M$.

Dla jednoznaczności zauważmy, że każda formuła z parametrami
z $A$ używa tylko skończonego podzbioru $F\subseteq A$.
W dowolnym nieskończonym rozszerzeniu elementarnym
wszystkie elementy spoza $F$ są zamienialne przez
permutację ustalającą $F$; jest ona automorfizmem struktury
w języku równości z nazwami $F$. Jeżeli dwa typy zawierające
wszystkie nierówności realizują się w różnych rozszerzeniach,
w każdym z nich możemy porównać realizację z tym samym
elementem $b\in M\setminus F$. Elementarność rozszerzeń
zapewnia tę samą prawdziwość danej formuły co dla $b$ w $M$.
Zatem oba typy rozstrzygają każdą formułę jednakowo.

Typ $p_a$ jest izolowany przez $x=c_a$. Jeśli $A$ jest
nieskończony, żadne otoczenie bazowe $p_\infty$ nie jest
singletonem: dla formuły z parametrami w $F$ każdy
$a\in A\setminus F$ spełnia ją, jeśli spełnia ją realizacja
$p_\infty$. Wtedy $S_1(A)$ jest uzwarceniem Aleksandrowa
dyskretnego $A$: otoczenia punktu $p_\infty$ zawierają
wszystkie poza skończenie wieloma punktami $p_a$,
a dopełnienia skończonych rodzin singletonów są otwarte.
Gdy $A$ jest skończony, także $p_\infty$ jest izolowany
przez koniunkcję wszystkich nierówności.
\end{przyklad}

\begin{przyklad}[Niezrealizowany typ przekroju wymiernego]
\label{ex:mod-type-cut}
Niech $\mathfrak M=(\Q,<)$, $A=\Q$ i $r\in\R\setminus\Q$.
Zbiór formuł
\[
p(x)=\{c_q<x:q\in\Q,\ q<r\}
\cup\{x<c_q:q\in\Q,\ r<q\}
\]
jest skończenie realizowalny w $\Q$ dzięki gęstości liczb
wymiernych w $\R$. Nie ma realizacji w $\Q$: dla dowolnego
$q_0\in\Q$ można znaleźć liczbę wymierną ściśle pomiędzy
$q_0$ i $r$, co daje formułę niespełnioną przy $q_0$.
Każde uzupełnienie $p$ do typu zupełnego nad $A$ jest
pomijane w $\mathfrak M_A$, ale realizowane w pewnym
rozszerzeniu elementarnym.
\end{przyklad}

\begin{przyklad}[Typ elementu przestępnego]
\label{ex:mod-type-transcendental}
Niech $K=\overline{\Q}$ będzie domknięciem algebraicznym
$\Q$ w języku pierścieni, z parametrami $A=\Q$.
Zbiór formuł $p(x)=\{f(x)\ne0:0\ne f\in\Q[X]\}$,
gdzie współczynniki są nazwanymi parametrami,
jest skończenie realizowalny w $K$: skończenie wiele
niezerowych wielomianów ma łącznie skończenie wiele pierwiastków,
a $K$ jest nieskończone. Nie jest realizowany w $K$, bo
każdy element $K$ jest algebraiczny nad $\Q$.
Twierdzenie \ref{thm:mod-type-elementary-extension} daje
rozszerzenie elementarne $K$ zawierające element przestępny
nad $\Q$.
\end{przyklad}

\section{Odwzorowania między przestrzeniami typów}
\begin{twierdzenie}[Ograniczenie parametrów]
\label{thm:mod-type-restriction}
Dla $A\subseteq B\subseteq M$ odwzorowanie
\[
r_{BA}:S_n(B)\longrightarrow S_n(A),\qquad
p\longmapsto p\cap F_n(\mathcal L(A))
\]
jest ciągłe i na.
\end{twierdzenie}
\begin{proof}
Ograniczenie typu jest zupełne w mniejszym języku
i skończenie realizowalne w $\mathfrak M_A$
na mocy twierdzenia \ref{thm:mod-type-elementary-extension}.
Jest więc elementem $S_n(A)$.
Dla $\varphi\in F_n(\mathcal L(A))$ przeciwobraz
zbioru bazowego $[\varphi]$ jest zbiorem bazowym
wyznaczonym przez tę samą formułę w $S_n(B)$.
To daje ciągłość. Dla $q\in S_n(A)$ wybieramy
jego realizację $\bar b$ w rozszerzeniu elementarnym
$\mathfrak N\succeq\mathfrak M$; wtedy
$\operatorname{tp}^{\mathfrak N}(\bar b/B)$ jest typem
nad $B$ (ponieważ $\mathfrak N_B\models T_B$), którego
ograniczenie jest równe $q$.
\end{proof}

\begin{uwaga}
Rozróżniamy pomijanie typu w ustalonym modelu od twierdzenia
o pomijaniu typów, które przy dodatkowych hipotezach
konstruuje model pomijający zadane typy. To drugie twierdzenie
wymaga osobnego dowodu. Następnym krokiem jest badanie modeli,
które realizują możliwie wiele typów, czyli modeli nasyconych.
\end{uwaga}
